\chapter{Latin Squares and Other Permutation Counts}
\label{ch:latin}

The queens are one of a family of counting problems on permutations in which an explicit constraint cuts an exponentially small fraction out of $N!$, and the family has a shared asymptotic theory, a shared set of methods, and a shared open question about constants. This chapter surveys the family: derangements, where everything is exact; Latin squares, where the leading constant is known and the lower order terms are open; the toroidal queens, where the count is a number theoretic function; and the general permanent of a zero one matrix, where the van der Waerden and Br\`egman bounds frame everything.

\section{Derangements and the rook polynomial}

A derangement is a permutation with no fixed point, and the number $D(N)$ of derangements is $N!$ times $\sum_{k=0}^{N} (-1)^k / k!$, which is $N!/e$ rounded to the nearest integer. The count is exact by inclusion exclusion, and the ratio $D(N)/N! \to 1/e$ is the simplest instance of a rare event in permutations that is not rare: a constant fraction. More generally, for a board $B$ of forbidden squares, the number of permutations avoiding $B$ is $\sum_k (-1)^k r_k(B) (N - k)!$ where $r_k$ is the number of ways to place $k$ non attacking rooks on $B$, the rook polynomial of Kaplansky and Riordan. For the queens the forbidden structure is not a fixed set of squares but a set of pairs, which is why inclusion exclusion does not close and the problem is $\sharpP$ hard.

\section{Latin squares}

A Latin square of order $n$ is an $n \times n$ array over $n$ symbols with each symbol once in each row and column; equivalently it is a sequence of $n$ permutations that are pairwise disjoint as sets of positions. The number $L(n)$ is known exactly only to $n = 11$, where it is about $7.8 \times 10^{47}$, and the leading asymptotics follow from the van der Waerden and Br\`egman inequalities on permanents, as presented by van Lint and Wilson, which give
\[
L(n)^{1/n^2} \sim \frac{n}{e^2} ,
\]
so that $\log L(n) = n^2 \log n - 2 n^2 + o(n^2)$. The constant $e^{-2}$, the analogue of Simkin's constant for the queens, is known; the lower order terms are not: van der Waerden's inequality and Br\`egman's inequality, applied row by row, pin $\log L(n)$ to $n^2 \log n - 2n^2 + o(n^2)$ and leave the terms of order $n \log n$ and $n$ open. Keevash's existence and counting theorems for designs, and the refinements of Kwan, Sah, Sawhney and Simkin on substructures in Latin squares, sharpen the error term and count Latin squares with prescribed structure, but the constant in the next order remains unidentified.

The method of the Latin square counts is the method of the queens: entropy upper bounds via the chain rule, applied row by row, and random greedy or absorption lower bounds. Simkin's queens paper is explicitly in this tradition and borrows the technology of the Latin square and design counts. The queens differ in that the constraint is between every pair of rows at once rather than between a row and the set of previous rows, which is what makes the queens a model with a nontrivial limit object, the queenon, where the Latin square limit is trivial.

\section{Toroidal queens}

Place the queens on a torus, so that the diagonals wrap around. The number of toroidal solutions $T(N)$ is zero unless $\gcd(N, 6) = 1$, by a theorem of P\'olya from 1918, and for such $N$ the linear permutations $x_i = a i + b \bmod N$ with $a$, $a - 1$ and $a + 1$ all coprime to $N$ are solutions, which already gives a count growing nearly quadratically in $N$; the exact values are known for small $N$ and grow as $((1 + o(1)) N e^{-\alpha_T})^N$ for a toroidal constant $\alpha_T \ge \alpha$, since every toroidal solution is a planar solution. The torus is strictly rarer, and its constant is known: Luria~\cite{Luria2017} proved the upper bound $T(N) \le ((1 + o(1)) N e^{-3})^N$, and Bowtell and Keevash~\cite{BowtellKeevash} proved the matching lower bound for $\gcd(N, 6) = 1$, so that $\alpha_T = 3 > \alpha \approx 1.942$. This is what the translation invariant variational problem on the torus suggests, with the uniform density as optimiser, and it is one of the few exact constants in the family.

The toroidal problem also has a level set structure that is more symmetric than the planar one, since the dihedral symmetry of the square is replaced by the affine group of the torus, and the swap chain on toroidal level sets should mix at least as well. We have not run it.

\section{The permanent and its bounds}

The permanent of a zero one $N \times N$ matrix $A$ counts the permutations supported on the ones of $A$, and derangements and permutations avoiding a board are permanents. The two classical inequalities are van der Waerden's, proved by Egorychev and Falikman in 1981, that a doubly stochastic matrix has permanent at least $N! / N^N$, and Br\`egman's of 1973, that a zero one matrix with row sums $r_i$ has permanent at most $\prod_i (r_i!)^{1/r_i}$. Together they bracket the number of permutations avoiding a given forbidden board to within a factor exponential in $N$, which is the precision of the leading order asymptotics for Latin squares and for every other problem of the family. The queens are not a single permanent, because the constraint is on pairs of positions rather than on a fixed board, and the standard expansion is an inclusion exclusion over sets of attacking pairs, which has no permanent structure.

Jerrum, Sinclair and Vigoda's FPRAS for the permanent, discussed in Chapter~\ref{ch:labs}, applies to every single permanent in the family and so gives a polynomial time approximation for derangements, for permutations avoiding any fixed board, and for the number of ways to extend a Latin rectangle by one row. It does not apply to the queens, whose inclusion exclusion over attacking pairs has exponentially many terms with alternating signs, and the alternation is what defeats the method: the FPRAS for a sum with cancellation is not a sum of FPRAS for the terms. This is, from the permanent side, the reason the queens are harder than derangements and the other single permanents of the family.

\section{The constants}

Every problem in the family has an asymptotic of the form $( (1 + o(1)) N c )^N$ or $( (1 + o(1)) n c )^{n^2}$ with an explicit $c$: $1/e$ for derangements, $1/e^2$ for Latin squares, $e^{-\alpha}$ for the queens. The first is exact by inclusion exclusion. The second is exact by van der Waerden and Br\`egman. The third is a variational constant computed numerically and not known in closed form, and the difference is that the queens have a nontrivial limit object. The family suggests that the right question about $\alpha$ is not its decimal expansion but whether it is the value of a recognisable functional at a recognisable density, in the way that $1/e^2$ is the value of the Br\`egman bound at the uniform doubly stochastic matrix. Simkin's queenon is the density; the functional is his; and the question whether $\alpha$ can be identified is the question whether the queenon can be. We record it among the open problems of Chapter~\ref{ch:open} under the finite size and rate function headings.
